% ══════════════════════════════════════════════════════════════════
%  FICHE D'EXERCICES — Chapitre 6 : Les composés carbonylés (aldéhydes et cétones)
% ══════════════════════════════════════════════════════════════════
\section{Exercices type Baccalauréat}
\textit{Données : $M(\ce{H})=1$, $M(\ce{C})=12$, $M(\ce{O})=16$ \si{\gram\per\mol}.}

\rubrique{Nomenclature, identification}

\exo{}\diff{1} Nommer les composés et préciser s'il s'agit d'un aldéhyde ou d'une
cétone : a)~$\ce{CH3-CHO}$ ; b)~$\ce{CH3-CO-CH3}$ ; c)~$\ce{CH3-CH2-CHO}$ ;
d)~$\ce{CH3-CO-CH2-CH3}$.

\exo{}\diff{1} Écrire les formules semi-développées des composés carbonylés isomères
de formule brute $\ce{C4H8O}$, les nommer et les classer (aldéhyde / cétone).

\rubrique{Tests caractéristiques}

\exo{}\diff{1} Citer trois tests permettant de mettre en évidence un aldéhyde.
Préciser, pour chacun, l'observation attendue.

\exo{}\diff{2} On dispose de deux flacons contenant l'un du propanal, l'autre de la
propanone. Décrire un protocole expérimental permettant de les distinguer, en
indiquant les réactifs, les observations et en écrivant l'équation du test
caractéristique de l'aldéhyde (réactif de Tollens).

\exo{}\diff{2} Un composé $A$ de formule $\ce{C3H6O}$ donne un précipité jaune avec
la 2,4-DNPH mais ne réagit ni avec la liqueur de Fehling ni avec le réactif de
Tollens. Identifier la famille et la formule de $A$, et le nommer.

\rubrique{Obtention, oxydation}

\exo{}\diff{1} De quel alcool obtient-on, par oxydation ménagée : a)~le propanal ;
b)~la butan-2-one ? Écrire les formules des alcools et préciser leur classe.

\exo{}\diff{2} L'oxydation d'un aldéhyde par une solution de permanganate de
potassium acidifiée donne un acide carboxylique. Écrire l'équation de l'oxydation de
l'éthanal (couple $\ce{MnO4-}/\ce{Mn^2+}$).

\rubrique{Détermination de formule}

\exo{}\diff{2} Un composé carbonylé saturé $A$ a une masse molaire
$M = \SI{72}{\gram\per\mol}$.
\begin{enumerate}[label=\arabic*)]
\item Déterminer sa formule brute (formule générale $\ce{C_nH_{2n}O}$).
\item Écrire les isomères aldéhydes et cétones, les nommer.
\item $A$ donne le test de Tollens positif. Conclure sur sa nature.
\end{enumerate}

\exo{}\diff{3} La combustion complète de \SI{1.45}{\gram} d'un composé carbonylé
saturé $A$ fournit \SI{3.30}{\gram} de $\ce{CO2}$ et \SI{1.35}{\gram} d'eau.
\begin{enumerate}[label=\arabic*)]
\item Montrer que $A$ a pour formule brute $\ce{C3H6O}$.
\item $A$ ne réduit pas la liqueur de Fehling. L'identifier et le nommer.
\item De quel alcool $A$ peut-il provenir par oxydation ménagée ? Écrire l'équation.
\end{enumerate}

\rubrique{Problèmes de synthèse — type Baccalauréat}

\sujetbac[d'après Baccalauréat C, Cameroun]
Trois flacons $A$, $B$ et $C$ contiennent, dans le désordre : du butanal, de la
butan-2-one et du butan-1-ol. On réalise les tests suivants :
\begin{itemize}
\item $A$ et $C$ donnent un précipité jaune avec la 2,4-DNPH ; $B$ non.
\item $A$ réduit la liqueur de Fehling ; $C$ non.
\end{itemize}
\begin{enumerate}[label=\arabic*)]
\item Identifier $A$, $B$ et $C$. Justifier à partir des tests.
\item Écrire les formules semi-développées et les noms des trois composés.
\item Écrire l'équation de l'oxydation ménagée de $B$ en $A$ (oxydant en défaut).
\item Le composé $C$ peut-il subir une oxydation ménagée ? Justifier.
\end{enumerate}

\sujetbac[Réaction d'addition]
La propanone réagit avec l'acide cyanhydrique $\ce{HCN}$ par addition sur la double
liaison $\ce{C=O}$.
\begin{enumerate}[label=\arabic*)]
\item Rappeler la nature de la liaison du groupe carbonyle et expliquer pourquoi elle
peut subir une addition.
\item Écrire l'équation de l'addition de $\ce{HCN}$ sur la propanone et donner la
formule semi-développée du produit.
\end{enumerate}
